\section*{शब्दकोश प्रयोग निर्देश --- How to Use This Dictionary}

\vspace{0.2em}
{\small\color{bpgray}A guide to the lexicographical methodology, star priority system, visual nest hierarchy, and dialectal tags used in this dictionary.}
\vspace{0.6em}

\begin{multicols}{2}
\raggedcolumns

% --- Section 1: Star Rating System ---
{\color{bpblue}\bfseries\small व्यावहारिक प्राथमिकता (Star Priority System)}
\vspace{0.2em}

\begin{description}[leftmargin=0.8em, labelwidth=!, itemsep=0.3em, font=\bfseries\color{bpbluedark}]
  \item[\starsThree \ Core Spoken] दैनिक बोलचाल के अति-महत्वपूर्ण शब्द। ये शब्द हर \starsThree \ प्रविष्टि में भोजपुरी प्रयोग उदाहरण वाक्य के साथ प्रस्तुत किए गए हैं।
  \item[\starsTwo \ Formal / Regional] औपचारिक, लिखित, समाचार या किसी विशिष्ट क्षेत्र (जैसे पश्चिमी या दक्षिणी) में प्रयुक्त शब्द।
  \item[\starsOne \ Literary / Archaic] लोकगीतों, प्राचीन साहित्य या तकनीकी/कृषि क्षेत्र में प्रयुक्त दुर्लभ शब्द।
\end{description}

\vspace{0.5em}

% --- Section 2: Visual Hierarchy ---
{\color{bpblue}\bfseries\small पद पदानुक्रम (The Nest System)}
\vspace{0.2em}

\begin{description}[leftmargin=0.8em, labelwidth=!, itemsep=0.3em, font=\bfseries\color{bpbluedark}]
  \item[मुख्य शब्द (Root)] **अंक** --- बड़े बोल्ड अक्षरों में मुख्य धातु/शब्द।
  \item[व्युत्पन्न पद (Derivative)] $\hookrightarrow$ *अंकगणित* --- एक स्तर भीतर व्युत्पन्न शब्द।
  \item[समास/यौगिक (Compound)] $\rightarrow$ **अंकित मूल्य** --- दो स्तर भीतर मिश्रित या यौगिक पद।
\end{description}

\columnbreak

% --- Section 3: Dialect Tagging ---
{\color{bpblue}\bfseries\small क्षेत्रीय बोली भेद (Dialect Tags)}
\vspace{0.2em}

\begin{description}[leftmargin=0.8em, labelwidth=!, itemsep=0.3em, font=\bfseries\color{bpbluedark}]
  \item[{[पूर्वी]}] पूर्वी क्षेत्र (सारण, चंपारण, बलिया)।
  \item[{[पच्छिमी]}] पश्चिमी क्षेत्र (गोरखपुर, देवरिया, बस्ती)।
  \item[{[दक्षिणी]}] दक्षिणी क्षेत्र (भोजपुर, बक्सर, रोहतास)।
  \item[{[ठेठ]}] सर्वमान्य लोक-भोजपुरी रूप।
\end{description}

\vspace{0.5em}

% --- Section 4: Sample Entry ---
{\color{bpblue}\bfseries\small नमूना प्रविष्टि (Sample Entry Layout)}
\vspace{0.2em}

\begin{tcolorbox}[colback=bpbg,colframe=bpblue,arc=3pt,boxrule=0.5pt,top=4pt,bottom=4pt,left=6pt,right=6pt]
{\bfseries\large अंकवार} \ \K{\ba\banus\bka\bva\bmaa\bra} \, \textit{(Añkwaar)} \ \textbf{[सं. स्त्री. | ठेठ]} \ \starsThree \\
{\small\textbf{हिन्दी:} आलिंगन, गोद, छाती से लगाना।} \\
{\small\textbf{भोजपुरी:} अँकवार / कोरवा / गले लगावल} \\
{\small\textbf{Eng:} Embrace, open-armed hug, lap.} \\
{\small\color{bpbluedark}\textbf{उदा.:} ``दउड़ के माई के अंकवार भर लिहलसि।''} \\
{\small\color{bpgray}\textbf{देखें:} अंक, अँकवारना}
\end{tcolorbox}

\end{multicols}
